\documentclass[11pt]{amsart}
\usepackage{amsmath,amssymb,amsthm,mathtools}
\usepackage[margin=1.05in]{geometry}
\usepackage{booktabs}
\usepackage[hypertexnames=false]{hyperref}
\hypersetup{hidelinks}
\usepackage[final,expansion=false]{microtype}
\setlength{\parskip}{3pt plus 1pt}
\emergencystretch=2em

\newtheorem{theorem}{Theorem}[section]
\newtheorem{proposition}[theorem]{Proposition}
\newtheorem{lemma}[theorem]{Lemma}
\newtheorem{corollary}[theorem]{Corollary}
\theoremstyle{definition}
\newtheorem{definition}[theorem]{Definition}
\theoremstyle{remark}
\newtheorem{remark}[theorem]{Remark}
\newtheorem{observation}[theorem]{Numerical observation}

\newcommand{\R}{\mathbb R}
\newcommand{\C}{\mathbb C}
\newcommand{\Rel}{\operatorname{Re}}
\newcommand{\Tr}{\operatorname{Tr}}
\newcommand{\supp}{\operatorname{supp}}
\newcommand{\ran}{\operatorname{ran}}
\newcommand{\IT}{I_T}
\newcommand{\PT}{\mathcal P_T}
\newcommand{\Fh}{\widehat F}
\newcommand{\cF}{c_F}
\newcommand{\NT}{\mathcal N_T}
\newcommand{\Aconst}{\mathsf A_T}
\newcommand{\PiT}{\Pi_T}
\newcommand{\lamA}{\lambda_{\mathrm A}}
\newcommand{\Son}{\mathrm{Son}}
\newcommand{\Pih}{\widehat P}

\begin{document}

\title[A semilocal Sonin space and a finite reduction]{A semilocal Sonin space and a finite reduction of window positivity:\\ notes on a construction of Connes and Consani}
\author{David Alejandro Trejo Pizzo}
\thanks{Independent researcher, Buenos Aires. \texttt{dtrejopizzo@gmail.com}. October 2026.}
\date{}

\begin{abstract}
Starting from the Sonin space and the function $\varepsilon$ of Connes and Consani, we record three
things about the positivity of Weil's quadratic form on a window of half-length $T$. First, the
remainder left by the archimedean Sonin trace has infinite negative index on primitive tests as soon as
$T>\frac12\log2$, so the ``one negative direction'' structure of their archimedean theorem stops exactly
at $\log2$. Second, the semilocal Fourier transform $F_S=EF_\infty E^{-1}$ for $S=\{\infty,2\}$ is a
unitary self-adjoint involution whose scattering phase relative to $F_\infty$ reproduces the prime comb
of $2$ exactly; the compressed operator $PF_SP$ is compact but not Hilbert--Schmidt, the semilocal Sonin
projection has an operator gap, and the semilocal analogue of the Connes--Consani trace formula holds with
a remainder density having explicit logarithmic singularities at the powers of $2$, so that the window
form is a positive trace form minus a Hilbert--Schmidt form at every $T$. Third, an archimedean coercivity
theorem with explicit constants on functions satisfying the two Tate conditions blockwise reduces window
positivity to the positivity of an explicit real symmetric matrix of size $2k-2$. None of this produces a
sign; it localizes where positivity lives.
\end{abstract}

\maketitle

\section{Introduction}

\subsection*{A word before the mathematics}
I have been working on this problem for more than ten years, as a hobby, outside academia. Of all the
routes I have tried, the one I came to love is yours: the Sonin space, the function $\varepsilon$, the
prolate operators and the semilocal adelic picture that you built with Caterina Consani and Henri
Moscovici. It has a shape; one can see where the primes enter and what the archimedean place does on its
own.

I want to be honest about where I stand. I think the positivity of the quadratic form on windows of
growing length needs a genuinely new mathematical idea, and I do not have it. What I do have is a body of
work built directly on your construction: an exact semilocal version of your trace formula, a precise
statement of where your archimedean mechanism stops, a few theorems on the archimedean form that I have
not found in the literature, and a finite reduction that says exactly in which subspace the difficulty
sits. These notes collect that work, with proofs sketched and every statement labelled by what it is: a
theorem, a numerical observation, or an open question. If any of it helps you move the route forward, I
would be glad; if not, I would still like you to have it.

\subsection*{The object}
Throughout, $T>0$, $\IT=(-T,T)$, and for $F\in L^2(\IT)$ we write $\Fh(\tau)=\int F(t)e^{i\tau t}dt$ and
$\cF(u)=\int F(t)\overline{F(t-u)}dt$ for the autocorrelation, supported in $[-2T,2T]$. Put
\[
m(\tau)=\Rel\psi\big(\tfrac14+\tfrac{i\tau}2\big)-\log\pi,\qquad v_n=\frac{2\Lambda(n)}{\sqrt n},\qquad
\NT=\{n\ge2:\ \log n<2T\},\qquad \Aconst=\sum_{n\in\NT}v_n,
\]
and define the \emph{window form of half-length $T$} and its two pieces,
\begin{equation}\label{eq:form}
\mathfrak a_T(F)=\Gamma(F)+\PiT(F),\qquad
\Gamma(F)=\frac1{2\pi}\int_\R m(\tau)|\Fh(\tau)|^2d\tau,\qquad
\PiT(F)=-\sum_{n\in\NT}v_n\Rel\cF(\log n),
\end{equation}
on the form domain $\{F:\int\log(2+|\tau|)|\Fh|^2<\infty\}$. The two Tate moments are
$A(F)=\int Fe^{t/2}$, $B(F)=\int Fe^{-t/2}$, the \emph{primitive subspace} is $\PT=\{A(F)=B(F)=0\}$, and
$\lamA(T)$ denotes the bottom of $\mathfrak a_T$ on the unit sphere of $\PT$. The complete form
$W_T=\mathfrak a_T+2\Rel(A\bar B)$ is the one the explicit formula writes as a sum over zeros; on $\PT$ the
two agree. In your notation, $g\in C_c^\infty(\R_+^*)$, $f=g*g^*$ and $F(t)=g(e^t)$, so that $f(e^x)=\cF(x)$,
$\widehat g(0)=\int F$, $\widehat g(\mp i/2)$ are the two Tate moments, and $W_\infty(g*g^*)=\Gamma(F)$ by
\cite[(152)--(154)]{CC2021}; the two conditions ``vanishing at $i/2$ and $0$'' of your Theorem~1 are, up to
the second moment, the conditions defining $\PT$. The target throughout is $\lamA(T)\ge0$, window by window.
It is known with a certificate for $T\le0.8$ \cite{Zhu2026} and $\lamA\ge2.5\cdot10^{-12}$ for
$T\le\frac12\log5=0.8047$ \cite[Thm.~log5A]{Paper47}, and $\lamA(T)\to0$ at least as fast as
$T^2(k/T)^{2k}$ for every fixed $k$ \cite[Thm.~margin]{Paper47}. The form $\mathfrak a_T$ is also a L\'evy
energy, $\Gamma(F)=C_0\|F\|^2+\int_0^\infty(\|F\|^2-\Rel\cF(u))K(u)du$ with
$K(u)=2e^{-u/2}/(1-e^{-2u})$ and $C_0=\psi(\frac14)-\log\pi=-5.3722$ \cite[Lemma~K]{Paper47}; this
is the form in which the blocks of \S\ref{sec:coer} are computed.

\subsection*{What we asked and what we found}
Your archimedean theorem \cite[Thm.~1]{CC2021} writes $W_\infty(g*g^*)$ as the Sonin trace
$\Tr(\vartheta(g)\mathbf S\vartheta(g)^*)\ge0$ plus a remainder which, on the window $|I|=\log2$, has exactly
one negative direction, controlled by $|\widehat g(0)|^2$. We asked what this split does beyond
$\log2$, when primes enter, and whether the semilocal Sonin spaces of \cite{CCquasi,CCM2023,CCM2024} can
absorb them. In short:
\begin{itemize}
\item[(i)] the ``one negative direction'' structure stops exactly at $T=\frac12\log2$: for every larger $T$ the
remainder has infinite negative index on $\PT$, whatever bounded archimedean term one subtracts
(Theorem~\ref{thm:nogo});
\item[(ii)] for $S=\{\infty,2\}$ the semilocal Fourier transform is a unitary involution whose scattering phase is
the prime comb of $2$; the semilocal Sonin space absorbs the comb exactly and leaves a Hilbert--Schmidt
remainder at every window, with an explicit singular structure (Theorems~\ref{thm:phase}--\ref{thm:84});
\item[(iii)] the archimedean form is coercive, with explicit constants, on functions that satisfy the two Tate
conditions on each of $k$ consecutive blocks; hence $\lamA(T)>0$ if and only if an explicit real symmetric
matrix of size $2k-2$ is positive definite (Theorems~\ref{thm:B}, \ref{thm:schur}).
\end{itemize}
None of this produces a sign. What it does is localize where positivity lives: in the semilocal picture, in a
single almost Sonin-orthogonal direction where the window form is a $94\%$ cancellation between two positive
quantities (\S\ref{sec:num}); in the archimedean picture, in the span of the piecewise Tate exponentials
$e^{\pm t/2}\mathbf 1_{B_j}$ (\S\ref{sec:coer}).

\emph{Labels.} Statements called Theorem, Proposition, Lemma or Corollary are proved; the proofs below are
sketches, and the full arguments, with the independent re-derivations that were made of each of them, are in
the author's notes. Numerical observations are stated as such; none of them is a certificate.

\section{The window form and the Sonin identity}\label{sec:sonin}

Let $\mathbf S$ be the orthogonal projection on the Sonin space $S(1,1)$ of \cite[Def.~4.4]{CC2021} and
$\varepsilon$ the real, even, continuous function of \cite[Thm.~4.7, (83)--(84)]{CC2021}, so that
$\Tr(\vartheta(f)\mathbf S)=W_\infty(f)+\int f(\rho^{-1})\varepsilon(\rho)d^*\rho$ for every
$f\in C_c^\infty(\R_+^*)$. For $F\in C_c^\infty(\IT)$ put
\[
\Son(F):=\Tr\big(\vartheta(g)\mathbf S\vartheta(g)^*\big)\ \ge0,\qquad
\langle C_\varepsilon F,F\rangle:=\int_{-2T}^{2T}\cF(x)\,\varepsilon(e^{|x|})\,dx ,
\]
where $C_\varepsilon$ is the operator on $L^2(\IT)$ with the continuous kernel $\varepsilon(e^{|x-y|})$; it is
Hilbert--Schmidt.

\begin{theorem}[the Sonin split]\label{thm:split}
For every $T>0$ and every $F\in C_c^\infty(\IT)$, with no moment condition,
$\mathfrak a_T(F)=\Son(F)-\langle C_\varepsilon F,F\rangle+\PiT(F)$.
\end{theorem}
\begin{proof}
This is (83) read in the additive variable: $W_\infty(f)=\Gamma(F)$ by (152)--(154), $\Tr(\vartheta(f)\mathbf S)=\Son(F)$ by
cyclicity ($\vartheta(g)$ bounded, $\mathbf S\vartheta(g)^*$ Hilbert--Schmidt), and
$\int f(\rho^{-1})\varepsilon(\rho)d^*\rho=\int\cF(x)\varepsilon(e^{|x|})dx$; the prime comb only sees $k\log p<2T$ by the
support of $\cF$.
\end{proof}

For $T\le\frac12\log2$ the comb is empty and the remainder $-C_\varepsilon$ is compact; your \S6 shows, in the
metric of the variable $\xi$ ($F=\frac12\xi-\xi'$), that it has exactly one negative direction at $T=\frac12\log2$:
$\lambda_{\max}(K_I)\approx1.05158$, $\lambda_2(K_I)\le0.7722$, with the negative eigenvector $\zeta$ overlapping
the constant by $\langle\zeta|\xi_0\rangle\approx0.94865$; and by Cauchy interlacing the same ``at most one''
holds for every $T<\frac12\log2$. The next statement says that this is the whole range of that mechanism.

\begin{theorem}[infinite negative index past $\frac12\log2$]\label{thm:nogo}
Let $T>\frac12\log2$, let $k\in L^\infty(-2T,2T)$ be real and even, and let $C_k$ be the operator with kernel
$k(x-y)$ on $L^2(\IT)$. For every $N\ge1$ there is an $N$-dimensional subspace of $C_c^\infty(\IT)\cap\PT$ on which
$-\langle C_kF,F\rangle+\PiT(F)<0$ for $F\neq0$. In particular $\mathfrak a_T-\Son$ has infinite negative index on
$\PT$, also on $\PT\cap\{\int F=0\}$, and no inequality
$\mathfrak a_T-\Son\ge-b|v\rangle\langle v|+c(I-|v\rangle\langle v|)$ with $c>0$, nor any with finitely many negative
directions, holds on $\PT$ for any inner product.
\end{theorem}
\begin{proof}
Fix $c_0\in C_c^\infty(-1,1)$ real, $c_w=c_0(\cdot/w)$ and $\beta=c_w''-\frac14c_w$; two integrations by parts give
$\int\beta e^{\pm t/2}=0$, so all translates of $\beta$ are in $\PT$. Put $a=\log2$ and
$h_j=\beta(\cdot-t_j)+\beta(\cdot-t_j-a)$, $j=1,\dots,N$, with the $t_j$ chosen in $(-T+w,T-a-w)$ so that all mutual
differences avoid the finite set $\{\pm\log n\pm\{0,a\}:n\in\NT\}$ by more than $4w$; this is possible because
$2T-a>0$. For $x=\sum c_jh_j$ the only non-vanishing value of $\cF$ at a contact is the inner one of each pair, at
$\log2$, so $\PiT(x)=-v_2\|\beta\|_2^2\sum|c_j|^2=-\sqrt2\log2\,\|\beta\|_2^2\sum|c_j|^2$, while
$|\langle C_kx,x\rangle|\le4N\|k\|_\infty\|\beta\|_1^2\sum|c_j|^2$ and $\|\beta\|_1^2/\|\beta\|_2^2=O(w)$. For small $w$
the form is $\le-\frac12\sqrt2\log2\,\|\beta\|_2^2\sum|c_j|^2$. The last assertion is a dimension count. With
$k=\varepsilon(e^{|\cdot|})$, bounded on $[0,2T]$ by continuity of $\varepsilon$, this is $\mathfrak a_T-\Son$.
\end{proof}

\begin{remark}
(1) The theorem uses neither the sign of $\mathfrak a_T$ nor the identity of the prime. It does not contradict
$\mathfrak a_T\ge0$: on the test pairs, $\Son$ grows like $w^{-3}\log(1/w)$ and the comb like $w^{-3}$. It says that the
archimedean Sonin trace, subtracted \emph{in full}, is the wrong positive part once a contact is present. (2) Ceding a
fraction restores finiteness: for every $\eta\in(0,1]$, $n_-(\mathfrak a_T-(1-\eta)\Son|_{L^2(\IT)})<\infty$, bounded by the
number of eigenvalues $\ge A'/(A'+L)$ of the time--band operator of half-length $T$ and bandwidth $s_*$, the least $s$
with $\eta m(s)\ge2A'$, $A'=\|C_\varepsilon\|+\Aconst$; since $m(s)=\log(s/2\pi)+O(s^{-2})$ this is of order
$Te^{2A'/\eta}$. The index is finite, but the finite arena it leaves is exponentially large in $1/\eta$.
\end{remark}

\section{A semilocal Sonin space for $S=\{\infty,2\}$}\label{sec:semilocal}

\subsection{The model and the transform}
We work in $H=L^2(\R_+,du)$, which is $L^2(\R)^{\mathrm{ev}}$ with your inner product $\frac12\int_\R=\int_0^\infty$
\cite[(16)]{CC2021}; $F_\infty$ is the operator with kernel $2\cos(2\pi ux)$, $P=\mathbf 1_{(0,\lambda]}(u)$ is the cut-off,
$(\vartheta(\rho)\xi)(u)=\rho^{-1/2}\xi(\rho^{-1}u)$ is the scaling action, and $\vartheta(f)=\int f(\rho)\vartheta(\rho)d^*\rho$.
In the logarithmic variable $t=\log u$ (unitarily, $\xi\mapsto e^{t/2}\xi(e^t)$) scalings become translations and we write
$D=-i\,d/dt$. Let $D_2\xi(u)=\xi(2u)$ and $E=\sum_{k\ge0}D_2^k$, so $E^{-1}=1-D_2$; in Mellin form on $\Rel s=\frac12$,
$E$ is the bounded multiplier $(1-2^{-s})^{-1}$. Following \cite[Prop.~4.1]{CCM2023}, where $F_S\circ\eta_S=\eta_S\circ F_\infty$
with $\eta_S=E$ in this model, the \emph{semilocal Fourier transform} is
\[
F_S:=EF_\infty E^{-1}.
\]
Put $\Pih^S=F_SPF_S$, $Q_1=1-P$, $Q_2=1-\Pih^S$, $A_S:=PF_SP$, and let $\Pi_S$ be the orthogonal projection on
$\ran Q_1\cap\ran Q_2$, the semilocal Sonin space (your $\ker u_{22}$ in \cite{CCquasi}). The level $\lambda$ is a free
parameter; the archimedean formula (83) is the case $\lambda=1$.

\begin{theorem}[scattering phase]\label{thm:phase}
In the logarithmic variable, $F_\infty=I\,m_\infty(D)$ with $I$ the reflection $t\mapsto-t$ and $|m_\infty|=1$ (the Tate
factor), and $F_S=F_\infty\,e^{i\varphi(D)}$ with
\[
\varphi'(t)=-\sum_{k\ge1}v_{2^k}\cos(kt\log2),\qquad v_{2^k}=2\log2\cdot2^{-k/2}.
\]
For finite $S$ with $F_S=E_SF_\infty E_S^{-1}$, $E_S=\prod_pE_p$, the phases add: $\varphi_S'=-\sum_{p\in S}\sum_kv_{p^k}\cos(kt\log p)$.
Consequently, for every $\lambda>0$ and every $h$ with $y\mapsto h(e^y)$ Schwartz, $\vartheta(h)(\Pih^S-\Pih_\infty)$ is trace class and
\begin{equation}\label{eq:Cprime}
\Tr\,\vartheta(h)\big(\Pih^S-\Pih_\infty\big)=\sum_{p\in S}\sum_{k\ge1}\log p\;p^{-k/2}\big(h(p^k)+h(p^{-k})\big),
\end{equation}
exactly, with no remainder and no dependence on $\lambda$.
\end{theorem}
\begin{proof}
With $q=2^{-1/2}$ and $\mathsf T$ the translation by $\log2$, $E^{-1}=1-q\mathsf T$, $F_\infty EF_\infty=E^*$, so
$F_S=F_\infty U$ with $U=E^*E^{-1}=(1-q\mathsf T)/(1-q\mathsf T^{-1})$, a quotient of conjugates: the symbol is
$e^{i\varphi(t)}$, $\varphi(t)=2\arg(1-qe^{it\log2})=-2\sum_kq^k\sin(kt\log2)/k$, and differentiating gives the comb. This $U$ is
the factor $\rho_2(z)=\gamma_2(z)/\gamma_2(1-z)$ of your quasi-inner function $u=\rho_\infty\prod_p\rho_p$
\cite[Lemma~3.1]{CCquasi}; we only verified the normalization. For \eqref{eq:Cprime} write $U=\sum_ju_j\mathsf T^{j}$ with
$\sum_j(1+|j|)|u_j|<\infty$ (geometric coefficients); then $UP_aU^*-P_a=\sum_{j,l}u_j\bar u_l(P_{a-j\log2}-P_a)\mathsf T^{j-l}$,
each term is trace class against $\widehat k(D)$ by Simon's $\ell^1(L^2)$ criterion \cite[Thm.~4.5]{Simon}, and the trace of the
series is $-\frac1{2\pi}\int\widehat k\,\varphi'$; conjugating by $F_\infty$ reverses $D$. Additivity in $S$ follows from the
identity $Y(\varphi_1+\varphi_2)=e^{i\varphi_1(D)}Y(\varphi_2)e^{-i\varphi_1(D)}+Y(\varphi_1)$ for
$Y(\varphi)=e^{i\varphi(D)}P_ae^{-i\varphi(D)}-P_a$ and invariance of the trace.
\end{proof}

The scattering phase of $F_S$ relative to $F_\infty$ is therefore the prime comb of the window form, coefficient by
coefficient, with $\Lambda$ appearing as a logarithmic derivative; \eqref{eq:Cprime} is the dual form of the classical
phase-shift trace formula and is the local term at $p$ of your semilocal trace formula. It concerns the \emph{linear}
object $1-\Pih^S$; the nonlinear part, the intersection $\ran Q_1\cap\ran Q_2$, is where the rest of this section lives.

\begin{proposition}[the compressed transform]\label{prop:AS}
\emph{(i)} $F_S$ is a real, unitary, self-adjoint involution and $F_S\vartheta(\rho)F_S=\vartheta(\rho^{-1})$; hence
$A_S=PF_SP$ is compact and self-adjoint, $P\Pih^SP=A_S^2$, and $\|A_S\|<1$ ($\pm1$ are not eigenvalues).
\emph{(ii)} $A_S$ is not Hilbert--Schmidt: $\sum_n\lambda_n(A_S)^2=\infty$. It belongs to the Schatten class
$\mathcal S_p$ for every $p>2$.
\end{proposition}
\begin{proof}
(i) $F_S$ is similar to the involution $F_\infty$ and unitary by Theorem~\ref{thm:phase}; $U$ commutes with the scalings.
Compactness is yours \cite[Thms.~2.1, 4.4]{CCquasi}; directly, the grouped kernel $\sum_{k\le K}(2c_k-c_{k-1})$,
$c_k=\cos(2\pi2^kux)$, gives Hilbert--Schmidt operators converging to $A_S$ in norm, with
$\|PD_2^kF_\infty E^{-1}P\|\le(1+q)q^k$. If $A_S\xi=\pm\xi$ then $F_S\xi=\pm\xi$ is supported in $(0,\lambda]$, so
$E^{-1}\xi$ and its cosine transform are both compactly supported, forcing $\xi=0$.
(ii) If $A_S$ had an $L^2$ kernel $\kappa$, pairing with $\chi c_m$, $\chi\in C_c^\infty((\delta,\lambda-\delta)^2)$, would
give $\langle\kappa,\chi c_m\rangle\to0$ by Riemann--Lebesgue, while the lacunary structure of the kernel (non-stationary phase
in every term $c_jc_m$, $j\ne m$) gives $\langle\kappa,\chi c_m\rangle=\frac12\iint\chi+O(2^{-m})\ne0$. For $p>2$,
$\|Pc_kP\|_{\mathcal S_p}\le\|Pc_kP\|^{1-2/p}\|Pc_kP\|_{\mathrm{HS}}^{2/p}\le(\tfrac12q^k)^{1-2/p}\lambda^{2/p}$ is summable.
\end{proof}

So the primes change the spectrum of the band operator genuinely: the similarity $A_S=EA_\infty E^{-1}$ is false
(the defect is $PE(1-P)F_\infty PE^{-1}P$, the mass outside the box folded back by the dyadic dilations), and the
superexponential Slepian decay of the archimedean eigenvalues is replaced by a square-divergent sequence. At
$\lambda=\frac12$ the first eigenvalues are $0.988$, $-0.507$ (numerical) against $0.885$, $-0.107$ for $A_\infty$;
at $\lambda=1$, $1-\lambda_1\approx2.6\cdot10^{-6}$ and $\lambda_2\approx-0.9927$ (numerical, converged to the digits shown),
so the gap of the next theorem is real but tiny at your level.

\subsection{Gap, trace class and the remainder density}
Let $\{\xi_n\}$ be a real orthonormal eigenbasis of $A_S$ in $\ran P$, $A_S\xi_n=\lambda_n\xi_n$, and for $\lambda_n\ne0$
put $\zeta_n=(1-P)F_S\xi_n/(1-\lambda_n^2)^{1/2}$ and $\tau_n=\lambda_n(1-\lambda_n^2)^{-1/2}$.

\begin{theorem}[operator gap and trace class]\label{thm:gap}
\emph{(i)} In Halmos' decomposition of the pair $(Q_1,Q_2)$ the angles satisfy $\cos^2\theta_n=\lambda_n^2$, the $\zeta_n$
are orthonormal, $Q_1Q_2Q_1=\Pi_S+\sum_n\lambda_n^2|\zeta_n\rangle\langle\zeta_n|$, and
$\operatorname{spec}(Q_1Q_2Q_1|_{\ran Q_1})\subset\{0\}\cup\{\lambda_n^2\}\cup\{1\}$ with $\sup\lambda_n^2=\|A_S\|^2<1$. Hence
$\Pi_S$ is the Riesz projection of $Q_1Q_2Q_1$ at the isolated point $1$, and
$\|(Q_1Q_2Q_1)^n-\Pi_S\|=\|A_S\|^{2n}$.
\emph{(ii)} For $g,f\in C_c^\infty(\R_+^*)$: $\vartheta(g)Q_1Q_2$ is Hilbert--Schmidt; $\vartheta(f)Q_1Q_2Q_1$ and
$\vartheta(f)\Pi_S$ are trace class; and $P\vartheta(f)P\,A_S^2$ is trace class, so that
$\sum_n\lambda_n^2\langle\xi_n|\vartheta(f)\xi_n\rangle$ converges absolutely although $\sum\lambda_n^2=\infty$.
\end{theorem}
\begin{proof}
(i) is the two-projection calculus \cite{Halmos1969} with $P\Pih^SP=A_S^2$ from Proposition~\ref{prop:AS}(i).
(ii) By Theorem~\ref{thm:phase}, after the reflection $I$, $Q_1Q_2=I\,\mathbf 1_{t<-a}M(D)\mathbf 1_{t>a}M(D)^*I$ with
$M=m_\infty e^{i\varphi}$ unimodular, $a=\log\lambda$: the semilocal case is literally the archimedean one with another
symbol. Then $\widehat{\tilde g}(D)\mathbf 1_{t<-a}M(D)\mathbf 1_{t>a}$ splits into a commutator of Hilbert--Schmidt norm
$\int|k|^2|x|dx$ and a Hankel block of norm $\int|K_g|^2(-2a-x)_+dx$, finite because $\widehat{\tilde g}\,m_\infty e^{i\varphi}$ is
Schwartz ($m_\infty$ has polynomially bounded derivatives by Stirling, $e^{i\varphi}\in C_b^\infty$). The trace-class
statements follow by $f=g_1*g_2$, your Lemma~D.1 for the commutators, and Dixmier--Malliavin;
$\vartheta(f)\Pi_S=\vartheta(f)Q_1Q_2Q_1\Pi_S$. For the last claim, $P\vartheta(f)(\Pih^S-\Pih_\infty)P=\Theta_{00}(A_S^2-A_\infty^2)+\Theta_{01}X_{10}$
with $\Theta_{01}=P\vartheta(f)(1-P)$ trace class and $A_\infty$ trace class; the left side is trace class by \eqref{eq:Cprime}.
\end{proof}

\begin{theorem}[structure of the remainder density]\label{thm:L2}
The series $E_S(f):=\sum_n\tau_n\big(\langle\xi_n|\vartheta(f)\zeta_n\rangle+\langle\zeta_n|\vartheta(f)\xi_n\rangle\big)$ converges
absolutely for every $f\in C_c^\infty(\R_+^*)$ and equals $\int f(\rho^{-1})\varepsilon_S(\rho)d^*\rho$ with
$\varepsilon_S(\rho^{-1})=\varepsilon_S(\rho)$ and, for $\rho\ge1$,
$\varepsilon_S(\rho)=\rho^{1/2}\int_{\lambda/\rho}^\lambda G(\rho x,x)dx$, $G=(1-P)F_SA_S(1-A_S^2)^{-1}P$. For every $R>1$ and
$1\le\rho\le R$,
\[
\varepsilon_S(\rho)=\sum_{0\le j\le\log_2R}\beta_j\,\log^+\frac1{|\rho-2^j|}+O_R(1),\qquad
\beta_j=2^{j/2}\big(1-2^{-j}\big)\frac{\lambda^2}{2\log2},
\]
and $\varepsilon_S(\rho)\sim\frac{\lambda^2}{2\log2}(\rho-1)\log\frac1{\rho-1}$ as $\rho\downarrow1$. Hence
$\varepsilon_S\in L^p_{\mathrm{loc}}(\R_+^*)$ for every $p<\infty$, $\varepsilon_S\notin L^\infty_{\mathrm{loc}}$, and the pointwise
series diverges to $+\infty$ at every $\rho=2^j$, $j\ge1$.
\end{theorem}
\begin{proof}
Absolute convergence: the two cross series are traces of $\Theta_{01}G$ and $\Theta_{10}G^*$ in an orthonormal basis.
Write $G=G_0(1+R)$ with $G_0=(1-P)F_SPF_SP$ and $R=A_S^2(1-A_S^2)^{-1}$, Hilbert--Schmidt by Proposition~\ref{prop:AS}(ii), so
the $R$-part has bounded density. The kernel of $G_0$ is a double lacunary sum of $S(2^ku\mp2^lz)$,
$S(w)=\sin(2\pi\lambda w)/(2\pi w)$, and the lemma $\sum_{k\ge k_0}S(2^kt)=\lambda\,\#\{k\ge k_0:2^k|t|\le1\}+O(1)$, uniform in
$t\neq0$, gives $G_0(\rho x,x)=\frac{\lambda}{2\log2}\sum_i\log^+\frac1{|\rho-2^i|}+O_R(1)$ uniformly in $x\in[\lambda/R,\lambda]$;
integrate over $[\lambda/\rho,\lambda]$. Divergence at $2^j$: the dyadic quasi-covariance $K_S(2^jx,y)=K_S(x,y)-B_j(x,y)$, $B_j$
entire, reduces the pointwise term at $\rho=2^j$ to $\|\mathbf 1_{[2^{-j}\lambda,\lambda]}A_S\|_{\mathrm{HS}}^2$, which is infinite
by the argument of Proposition~\ref{prop:AS}(ii) localized to an interval. The numerical slopes of the three first
singularities, $0.128$, $0.272$, $0.448$ at $\lambda=\frac12$, agree with $\beta_1,\beta_2,\beta_3$ to one percent.
\end{proof}

The non-Hilbert--Schmidt character of $A_S$ thus shows up in exactly one place: the term $m=0$ of the Neumann series of
$(1-A_S^2)^{-1}$, whose density has positive logarithmic peaks at the powers of $2$, of weight growing like $2^{j/2}$. At
$\lambda=1$, $\beta_1=0.510$. There is no atom at $\rho=2$ in $\varepsilon_S-\varepsilon_\infty$: the atom is the comb, which
the next theorem separates.

\begin{theorem}[the semilocal trace formula]\label{thm:84}
Let $S=\{\infty,2\}$, $f\in C_c^\infty(\R_+^*)$ and $\mathrm{comb}_2(f)=\sum_{k\ge1}\log2\cdot2^{-k/2}\big(f(2^k)+f(2^{-k})\big)$.
\emph{(i)} For every $\lambda>0$,
$\Tr\,\vartheta(f)\Pi_S-E_S(f)=\Tr\,\vartheta(f)\Pi_\infty-E_\infty(f)-\mathrm{comb}_2(f)$.
\emph{(ii)} For $0<\lambda\le1$, the level dictionary
$\Tr\,\vartheta(f)\Pi_\infty^{(\lambda)}=W_\infty(f)+2\log\tfrac1\lambda\,f(1)+\int f(\rho^{-1})\varepsilon_\lambda(\rho)d^*\rho$ holds,
with $\varepsilon_1=\varepsilon$, and therefore
\[
\Tr\,\vartheta(f)\Pi_S=W_\infty(f)-\mathrm{comb}_2(f)+2\log\tfrac1\lambda\,f(1)+\int f(\rho^{-1})\varepsilon_S(\rho)\,d^*\rho .
\]
\end{theorem}
\begin{proof}
(i) For $f$ real with $f(\rho^{-1})=f(\rho)$ one has $[F,\vartheta(f)]=0$ for $F\in\{F_\infty,F_S\}$, and expanding
$\langle\xi_n|F\vartheta(f)F\xi_n\rangle$ with $F\xi_n=\lambda_n\xi_n+(1-\lambda_n^2)^{1/2}\zeta_n$ gives, for each $n$ and without
summing, $\langle\zeta_n|\vartheta(f)\zeta_n\rangle=\langle\xi_n|\vartheta(f)\xi_n\rangle-\tau_n(c_n+c_n')$, your (87). By
Theorem~\ref{thm:gap}(ii) the four series (two diagonal, two cross) converge absolutely, so your cancellation (89)--(90) can be
performed term by term: the diagonal $\lambda_n^2$ terms cancel $n$ by $n$ and
$\Tr\vartheta(f)Q_1Q_2Q_1-\Tr\vartheta(f)\Pi_F=E_F(f)+\sum\lambda_n^2\langle\zeta_n|\vartheta(f)\zeta_n\rangle$-type bookkeeping closes.
Subtracting the two cases, $Q_1Q_2^SQ_1-Q_1Q_2^\infty Q_1=-(1-P)(\Pih^S-\Pih_\infty)(1-P)$, whose trace against $\vartheta(f)$ is
computed blockwise and equals $-\mathrm{comb}_2(f)$ by \eqref{eq:Cprime}. Antisymmetric $f$ give zero on both sides since
$F\Pi_F=\Pi_FF$. (ii) $S(\lambda,\lambda)$ is unitarily $S(1,\lambda^2)$ by a scaling commuting with $\vartheta(f)$, and the
passage from $S(1,b)$ to $S(1,1)$, $b=\lambda^2\le1$, adds $\log\frac1b\,f(1)$ plus an explicit continuous density; combine with
(83) at $\lambda=1$.
\end{proof}

The term $2\log\frac1\lambda\,f(1)=2\log\frac1\lambda\,\|F\|^2$ was, for us, the difference between a consistent and an
inconsistent chain: at $\lambda=\frac12$ it is $1.386\|F\|^2$, and without it the semilocal Sonin trace came out negative on the
near-minimizer of the window form. With it, every check we could make passes, including the monotonicity
$\Son^{(1/2)}-\Son^{(1)}\succeq0$ (numerical margin $1.20$ at $T=\frac12\log3$).

\begin{corollary}[the window form at the first prime]\label{cor:split}
Let $\lambda=1$, $S=\{\infty,2\}$. For every $T>0$ and $F\in C_c^\infty(\IT)$,
\[
\Gamma(F)-\sum_{k:\,2^k\in\NT}v_{2^k}\Rel\cF(k\log2)=\Son_{\Pi_S}(F)-\langle C_{\varepsilon_S}F,F\rangle,\qquad
\Son_{\Pi_S}(F):=\Tr\big(\vartheta(g)\Pi_S\vartheta(g)^*\big)\ge0,
\]
and $C_{\varepsilon_S}$, the operator on $L^2(\IT)$ with kernel $\varepsilon_S(e^{|x-y|})$, is Hilbert--Schmidt. In particular, for
$T<\frac12\log3$, $\mathfrak a_T=\Son_{\Pi_S}-C_{\varepsilon_S}$ on $\PT$: the comb of $2$ is absorbed exactly, and the remainder
is compact, in contrast with Theorem~\ref{thm:nogo}. For $T\ge\frac12\log3$ the obstruction of Theorem~\ref{thm:nogo} returns
with the contact $3$, and the same chain for $S=\{\infty,2,3\}$ absorbs it, with $\varepsilon_S$ now carrying $\log^2$
singularities at the $S$-units $2^a3^b$ cut off at scale $1/b$ for $q=a/b$ in lowest terms.
\end{corollary}

\begin{remark}[attribution and what is not claimed]
From \cite{CCquasi,CCM2023,CCM2024}: $F_S$, the quasi-inner function and the factor $\rho_2$, the compactness of the Hankel
block, the semilocal Sonin space, and the moment and Jacobi structure for $S=\{p,\infty\}$. From \cite{CC2021}: the
archimedean (83)--(90), whose algebra Theorem~\ref{thm:84} transports. What we believe is new: Proposition~\ref{prop:AS}(ii),
Theorem~\ref{thm:gap} as stated (in particular the absolute convergence of the diagonal series), Theorem~\ref{thm:L2} with the
constants $\beta_j$, Theorem~\ref{thm:84} and Corollary~\ref{cor:split}. Nothing here decides a sign: the corollary writes
$\mathfrak a_T$ as a positive trace form minus a compact one, and the singular part of $\varepsilon_S$ is positive, so it
\emph{lowers} $\mathfrak a_T-\Son_{\Pi_S}$. For growing $S=S(T)=\{\infty\}\cup\{p<e^{2T}\}$ the constants are not uniform: a
heuristic count suggests that $\varepsilon_{S(T)}$ develops singularities of type $1/|\rho-n|$ at every integer $n<e^{2T}$, with
total $L^2$ mass growing exponentially in $T$. This is a conjecture, not a theorem, and it is the point at which the
semilocal remainder stops being ``archimedean plus a small perturbation''.\footnote{In a related but different
Legendre formulation, third-party interval certificates reach support $\log7$ and report that a semilocal form with the
prime $2$ alone loses positivity near $L\approx1.1147$, the odd sector failing first; we have not compared the two
formulations and do not claim that this says anything about $\Son_{\Pi_S}$.}
\end{remark}

\section{What the identity shows: the invariant $\rho^*(T)$ and the fragile direction}\label{sec:num}

Everything in this section is numerical (Galerkin on $\PT$ with $M\le32$ modes, the metric of $L^2(\IT)$, level $\lambda=1$,
$S=\{\infty,2\}$). The pipeline reproduces your constants: $\varepsilon'(1^+)=22.9965$, and at $T=\frac12\log2$, in the
$\xi$-metric, $\lambda_{\max}(K_I)=1.0504$ against your $1.05158\pm0.0012$, $\lambda_2=0.684$ against $0.6865$, and
$\langle\xi_0|\zeta\rangle=0.9488$ against $0.94865$.

By Corollary~\ref{cor:split}, $\mathfrak a_T=\Son_{\Pi_S}-C$ on $\PT$ for $T<\frac12\log3$, with $C=\langle C_{\varepsilon_S}\cdot,\cdot\rangle$
and $\Son_{\Pi_S}\succ0$ on $\PT$ (its minimum is $0.0145$ at $T=\frac12\log3$). Put
\[
\rho^*(T):=\sup_{F\in\PT}\frac{C(F)}{\Son_{\Pi_S}(F)},\qquad\text{so that}\qquad \mathfrak a_T\succ0\text{ on }\PT\iff\rho^*(T)<1 .
\]

\begin{observation}\label{obs:rho}
\[
\begin{array}{lcccccc}\toprule
T & \tfrac12\log2 & 0.35 & 0.40 & 0.45 & 0.50 & \tfrac12\log3\\ \midrule
\rho^*(T) & 0.001 & 0.003 & 0.055 & 0.245 & 0.644 & 0.944\\
C(u_T) & 0.003 & 0.006 & 0.021 & 0.026 & 0.023 & 0.017\\
\Son_{\Pi_S}(u_T) & 2.82 & 2.29 & 0.374 & 0.106 & 0.036 & 0.018\\ \bottomrule
\end{array}
\]
Here $u_T$ is the maximizer. The rise of $\rho^*$ from $0$ to $0.94$ is not driven by the numerator: $C(u_T)$ stays in
$[0.014,0.026]$ from $T=0.40$ on. It is driven by the collapse of the Sonin mass of the maximizer, which aligns with the
direction $s_0(T)$ of least Sonin mass in $\PT$ (overlap $0.998$ at $\frac12\log3$). At $T=\frac12\log3$ the maximizer is the
minimizer $h^*$ of $\mathfrak a_T$ itself: $\Son_{\Pi_S}(h^*)=0.0174$, $C(h^*)=0.0164$, $\mathfrak a_T(h^*)=9.9\cdot10^{-4}$. The
positivity of the window form at its weakest point is a $94\%$ cancellation between two positive quantities of the same size,
not the dominance of the Sonin trace over the remainder. The remainder does have genuinely negative directions
($-0.083$, $-0.040$), but there $\Son_{\Pi_S}$ wins by a factor $2.5$; and $h^*$ is orthogonal to them.
\end{observation}

\begin{observation}[the fragile direction is archimedean]\label{obs:s0}
The direction $s_0(T)$ of least Sonin mass for $\Pi_S$ coincides with the one for $\Pi_\infty$ (overlap $\ge0.989$ on the whole
range) and its minimal mass is the same ($0.0140$ archimedean, $0.0147$ semilocal, at $\frac12\log3$); it is even, with two sign
changes, a ``cosine with negative shoulders'', with $65\%$ of its mass in $|t|<T/2$ and a Fourier peak drifting like
$\approx4/T$ ($11.05$ at $\frac12\log2$, $7.60$ at $\frac12\log3$). It is not your $\zeta$ (overlap $0.03$): $\zeta$ maximizes the
remainder, $s_0$ minimizes the Sonin mass. The prime acts on $s_0$ only through $C$: $C_\infty(s_0)=-0.140$ becomes
$C_S(s_0)=+0.010$, and the balance in that direction is $\Delta C\approx0.135$ against $\mathfrak a^\infty_T(s_0)\approx0.140$
(the archimedean form in $s_0$); over all of $\PT$ the relative loss $\sup(\mathfrak a^\infty_T-\mathfrak a_T)/\mathfrak a^\infty_T$ is
$0.9936$. The mechanism is visible: the Fourier peak of $s_0$ sits where $\cos(\tau\log2)>0$ ($\cos(7.6\log2)=0.527$), so
$\Rel c_{s_0}(\log2)=+0.145>0$ and the comb costs $-v_2\cdot0.145=-0.142$, almost all of the archimedean surplus $0.153$.
\end{observation}

Two things we tried and that do not work. The truncated symbol $m+\widehat\varepsilon_T$ of $\Son_{\Pi_\infty}$ on the window is
not pointwise nonnegative ($-2.36$ at $\tau=0$ for $T=\frac12\log3$, with sign changes at $0.63$, $3.05$, $5.82$, $8.66$,
$11.14$): the positivity of the Sonin form on $\PT$ is a property of the compressed operator, and the small Sonin mass of $s_0$
is itself a cancellation ($+0.066-0.054$) between regions of either sign. And the count of negative directions of $-C_{\varepsilon_S}$
on $\PT$ is not a robust invariant: it grows from $2$ to $6$ under refinement because of the logarithmic peak at $\log2$, whose
displacement form has an oscillating symbol $\sim\cos(\tau\log2)/|\tau|$; only the two extreme eigenvalues are stable. What is
robust is $\rho^*(T)$ and the value of the two forms at $h^*$.

\section{Archimedean coercivity and the finite reduction}\label{sec:coer}

This part uses only the L\'evy form of $\Gamma$ and the elementary bounds
$\log\frac\tau{2\pi}-\frac1\tau-\frac1{2\tau^2}\le m(\tau)$ and $m$ increasing on $[0,\infty)$ with $m(0)=C_0$
\cite[Lemma~mbounds]{Paper47}. No Sonin space enters.

\begin{definition}
Let $k\ge1$, $\tau=T/k$, and let $B_1,\dots,B_k$ be the $k$ consecutive intervals of length $2\tau$ partitioning $\IT$. Put
\[
\mathcal V_k=\Big\{F\in L^2(\IT):\int_{B_j}Fe^{t/2}=\int_{B_j}Fe^{-t/2}=0\ (1\le j\le k)\Big\},\qquad E_k=\PT\ominus\mathcal V_k .
\]
\end{definition}
Then $\mathcal V_k\subset\PT$ is closed, $E_k=\operatorname{span}\{e^{\pm t/2}\mathbf 1_{B_j}\}\cap\PT$ has dimension exactly
$2k-2$, and $\PT=\mathcal V_k\oplus E_k$ orthogonally. Two structural facts about the blocks: the infimum of $\Gamma$ over
functions supported in one block of length $2\tau\le\log2$ with the two Tate conditions on that block is $\lamA(\tau)$, the bottom
of a prime-free window of half-length $\tau$, independently of the position of the block (translation invariance of the
autocorrelation); and for $F,G$ supported on disjoint blocks, $\Gamma(F+G)-\Gamma(F)-\Gamma(G)=-\iint F(x)G(y)K(y-x)$, in which
the rank-one leading term $2e^{-u/2}=2e^{-y/2}e^{x/2}$ of $K$ is killed by either Tate condition, leaving the reduced kernel
$K_1(u)=K(u)-2e^{-u/2}=2e^{-5u/2}/(1-e^{-2u})$, $0<K_1(u)\le u^{-1}$. We have not found this block structure in the literature.

\begin{theorem}[the short window]\label{thm:A}
For $0<\tau\le0.1$, $\lamA(\tau)\ \ge\ \log\frac1\tau-1.209-\tau$.
\end{theorem}
\begin{proof}
For $F\in\mathcal P_\tau$, $\|F\|=1$, the L\'evy form gives $\Gamma(F)=C_0+\int_{2\tau}^\infty K+\int_0^{2\tau}(1-\Rel\cF)K$; the tail is
$\log\coth\frac\tau2+\pi-2\arctan e^\tau\ge\log\frac2\tau+\frac\pi2-\tau$, and $K(u)\ge u^{-1}$ on $(0,2]$. Scaling to
$G(y)=\tau^{1/2}F(\tau y)$ on $[-1,1]$, the local integral is $\ge\int_0^2\frac{1-\Rel c_G(v)}vdv=\frac1{2\pi}\int w|\widehat G|^2$ with
$w(s)=\log(2|s|)+\gamma-\operatorname{Ci}(2|s|)$, and $w(s)\ge w(R)\min\{1,s^2/R^2\}$ for $0<R\le\pi$. The two Tate conditions say
$\widehat G(\mp i\tau/2)=0$, so $\widehat G=(s^2+\tau^2/4)\widehat h$ with $\supp h\subset[-1,1]$ and $h''\in L^2$; then
$1=\|G\|^2\ge\|h''\|^2\ge\frac{\pi^4}4\|h\|^2$ (double Poincar\'e inequality on $[-1,1]$), which bounds the low-frequency mass of
$\widehat G$ by $\frac{16(R^2+\tau^2/4)^3}{27\pi^4R^2}$. With $R=2$, $w(2)=2.1045$ and $C_0+\log2+\frac\pi2=-3.108$, the
constant follows.
\end{proof}
The true value is $\lamA(\tau)=\log\frac1\tau-0.510\pm0.001$ on $[0.02,0.1]$ (numerical, Galerkin). The qualitative fact, that
the archimedean form is positive in finite codimension on short windows, is Yoshida's Lemma~3 \cite{Yoshida1992}
(orthogonality to $N$ Fourier modes, no rate) and your Theorem~3.6 and Corollary~3.8 in \cite{CC2021} (the archimedean part is
$-2\cdot\mathrm{Id}$ plus Hilbert--Schmidt; positivity for $|I|\le2\cdot0.1404$); the leading term $\log\frac1a$ is observed
numerically for Suzuki's prime-free operator in \cite{Kim2026}. What seems new is the explicit rate with constants under
the two Tate constraints.

\begin{theorem}[coercivity on $\mathcal V_k$]\label{thm:B}
Let $\tau=T/k\le1$. For $F\in\mathcal V_k$ and $R>0$, $\|P_RF\|^2\le\delta(R\tau)\|F\|^2$ with
$\delta(\theta)=\frac{\cosh^2\tau}6(\theta^2+\frac{\tau^2}4)^2$, where $P_R$ is the Fourier projection on $|s|\le R$.
Consequently, for every $F\in\mathcal V_k$ in the form domain,
\[
\Gamma(F)\ \ge\ \Big(\log\frac1\tau-C_B(\tau)\Big)\|F\|^2,
\]
with $C_B(\tau)$ in closed form, $C_B(\tau)\to\frac14-\log\frac{6^{1/4}}{2\pi}=1.6399$ as $\tau\to0$, and
$C_B(0.02)=1.657$, $C_B(0.0619)=1.696$, $C_B(0.1)=1.733$.
\end{theorem}
\begin{proof}
On a block with centre $c$, subtract from $G=P_RF$ the element $u$ of $\operatorname{span}\{e^{\pm t/2}\}$ with $u(c)=G(c)$,
$u'(c)=G'(c)$; the remainder $r$ satisfies $r''-\frac r4=G''-\frac G4$ with zero Cauchy data at $c$, whence
$\int_Br^2\le\cosh^2\tau\frac{\tau^4}6\|G''-\frac G4\|^2_{L^2(B)}$. Since $F\mathbf 1_B\perp u$, $\langle F,G\rangle=\langle F,r\rangle$;
summing over blocks and using $\|G''-\frac G4\|\le(R^2+\frac14)\|G\|$ gives the leakage bound. Then
$\Gamma(F)/\|F\|^2=\int m\,d\mu=C_0+\int_0^\infty m'(\rho)\mu(|s|\ge\rho)d\rho\ge C_0+\int_0^{\theta_*/\tau}m'(\rho)(1-\delta(\rho\tau))d\rho$,
and an integration by parts with the lower bound for $m$ gives the closed form.
\end{proof}
Two momentsA of Tate per block make $\mathcal V_k$ a high-pass subspace, and $m\sim\log|\tau|$ does the rest. The theorem
deliberately does not separate the block-diagonal part from the coupling: on Tate-orthogonal functions on two adjacent
blocks the bilinear form $\iint F_1F_2K_1(|x-y|)$ has norm \emph{exactly} $\pi$ (Hilbert's inequality for $\le$; a Weyl sequence
$x'^{-1/2}\mathbf 1_{[2\tau e^{-M},2\tau]}$ in the distance to the common endpoint for $\ge$, the Tate projection costing only
$O(\sqrt M)$), so a bound of the form ``diagonal minus coupling norm'' would need $\log\frac1\tau>\pi+\|\PiT\|$, i.e.\ $k\approx160$
at $T=\frac12\log5$. In the Galerkin minimizers the coupling costs only $0.17$: the maximizers of the coupling norm are not
the minimizers of the form.

\begin{proposition}[the prime form]\label{prop:prime}
$|\PiT(F)|\le\Aconst\|F\|^2$. More precisely $\|\PiT\|\le\sum_{n\in\NT}v_n\cos\frac\pi{L_n+1}$ with $L_n=\lceil2T/\log n\rceil$,
which is $\Aconst/2$ for $T\le\log2$ ($=0.4901$ for $T\le\frac12\log3$) and $\frac{v_2}{\sqrt2}+\frac{v_3+v_4}2=1.6740$ for
$\log2<T\le\frac12\log5$.
\end{proposition}
\begin{proof}
$\Rel\cF(\ell)=\frac12\langle(S_\ell+S_\ell^*)F,F\rangle$ with $S_\ell$ the shift by $\ell$ compressed to $\IT$; fibred over
$r\in[0,\ell)$, $S_\ell+S_\ell^*$ is the adjacency matrix of a path with at most $L$ vertices, of norm $2\cos\frac\pi{L+1}$.
\end{proof}

\begin{corollary}[coercivity of the window form on $\mathcal V_k$]\label{cor:L1}
$c(T,k):=\inf\{\mathfrak a_T(F):F\in\mathcal V_k,\|F\|=1\}\ \ge\ \log\frac kT-C_B(\frac Tk)-\|\PiT\|$. For instance
$c(\frac12\log3,8)\ge0.486$ ($\dim E_8=14$), $c(\frac12\log5,25)\ge0.094$ ($\dim E_{25}=48$), $c(\frac12\log5,30)\ge0.281$.
For every $T$ there is $k_0(T)\le Te^{\Aconst+C_B}$ with $c(T,k)>0$ for $k\ge k_0(T)$.
\end{corollary}

\begin{theorem}[reduction to the piecewise Tate exponentials]\label{thm:schur}
Fix $T$ and $k$ with $c=c(T,k)>0$. Then $\mathfrak a_T$ is an inner product on $\mathcal V=\mathcal V_k\cap(\text{form domain})$
for which $\mathcal V$ is complete, for every $e\in E_k$ there is a unique $Ke\in\mathcal V$ with
$\mathfrak a_T(Ke,v)=\mathfrak a_T(e,v)$ for all $v\in\mathcal V$, and
\[
\mathfrak a_T(v+e)=\mathfrak a_T(v+Ke)+\mathfrak s_{T,k}(e),\qquad
\mathfrak s_{T,k}(e)=\mathfrak a_T(e)-\sup_{u\in\mathcal V}\big[2\Rel\mathfrak a_T(e,u)-\mathfrak a_T(u)\big].
\]
If $S_{T,k}$ is the matrix of $\mathfrak s_{T,k}$ in an orthonormal basis of $E_k$ and $\lambda_S$ its least eigenvalue, then
$\min\{c,\lambda_S\}/(1+\|K\|)^2\le\lamA(T)\le\min\{c,\lambda_S\}$. Hence $\lamA(T)>0$ if and only if the real symmetric
$(2k-2)\times(2k-2)$ matrix $S_{T,k}$ is positive definite.
\end{theorem}
\begin{proof}
Completeness: on $\mathcal V$, $\mathfrak a_T\ge c\|\cdot\|^2$ and $\mathfrak a_T+(\Aconst-C_0+1)\|\cdot\|^2$ is equivalent to the
weighted norm $\int\log(2+|s|)|\widehat v|^2$. The functional $v\mapsto\mathfrak a_T(e,v)$ is bounded by Cauchy--Schwarz for the
nonnegative form $\mathfrak a_T+(\Aconst-C_0)\langle\cdot,\cdot\rangle$, so $Ke$ exists by Riesz. The identity is the completion of
the square, and the sandwich follows from $\|v\|^2+\|e\|^2\le(1+\|K\|)^2(\|v+Ke\|^2+\|e\|^2)$.
\end{proof}

\begin{remark}[what the reduction is and is not]
It is an \emph{equivalence}, in an explicit basis of explicit dimension, with the constant $c$ coming from a prime-free
archimedean statement. It is not a gain in precision, and it is not the first finite reduction: sufficient finite reductions in
Legendre bases exist \cite[Thm.~1.1]{Zhu2026}, \cite{Groskin2026}, and the growth $k_0(T)\approx Te^{Ce^T}$ is of the same
doubly exponential order as Zhu's barrier \cite[Thm.~1.4]{Zhu2026}. By the sandwich, $S_{T,k}$ inherits the decay of the margin
in full: at $T=\frac12\log5$, $\lambda_S\approx\lamA\approx3\cdot10^{-12}$ with entries of size $O(1)$ (numerical). Its entries
are not in closed form, since the correction $\mathfrak a_T(Ke)$ is a supremum over $\mathcal V$; Galerkin gives upper bounds for
$S_{T,k}$, and a certified lower bound needs a near-maximizer together with the dual residual divided by $c$, which is
legitimate because $c>0$ is proved. The Galerkin near-minimizers of $\mathfrak a_T$ on $\PT$ have $99\%$ of their mass in $E_k$
for $k\ge5$; this is approximation theory (piecewise Tate exponentials approximate smooth functions to $O(\tau^2)$), not
arithmetic. The orbit blocks of the first contact ($k=3$, $\dim E=4$) do not give a $c$ uniform in $T$: it falls to $0.013$ at
$T=0.69$ (numerical). Related finite-section structure appears in \cite{CvS2025} (simple least eigenvalue with even
eigenfunction) and in Suzuki's screw-function operators on $[-a,a]$ \cite{Suzuki2023,Suzuki2026}; neither uses the Tate
conditions blockwise.
\end{remark}

\section{What we do not know}\label{sec:open}

The honest summary is that the two halves of this note localize the same difficulty from two sides without touching it.
On the semilocal side, the identity $\mathfrak a_T=\Son_{\Pi_S}-C_{\varepsilon_S}$ is exact and its remainder is compact, but at the
weakest direction of the window the two terms cancel to $6\%$, the remainder density has growing positive logarithmic peaks,
and we do not know a lower bound for the Sonin mass $\|\Pi_S\vartheta(g)\|^2_{\mathrm{HS}}$ in terms of $F$ that would see the
comb. On the archimedean side, everything is reduced to the positivity of $S_{T,k}$, whose entries are variational, and whose
least eigenvalue is already $10^{-12}$ at $T=0.805$. Beyond $T\approx0.85$ the least eigenvalue of the window form is at the
floor of double precision in every basis we used; we have no evidence of a sign there in either direction, and we have not tried
to go further than your range and Zhu's.

The questions we would most like to ask you. First, whether the semilocal Sonin space has, at your level $\lambda=1$, a
quantitative version of the operator gap of Theorem~\ref{thm:gap}: the gap $1-\|A_S\|^2$ is $5\cdot10^{-6}$ for $S=\{\infty,2\}$,
and we do not know how it behaves in $S$. Second, whether the moment structure of \cite{CCM2024} gives the density
$\varepsilon_S$, or at least its singular coefficients $\beta_j$, in closed form; our derivation goes through the lacunary kernel
and a Neumann series. Third, whether a differential operator commuting with $P$ and $\Pih^S$ exists at all: the candidate
$-\mathfrak s^2+\lambda^2N_S$ of \cite{CCM2023} does not commute with $P_\lambda^S$ for $S=\{\infty,2\}$ (the commutator against
$e^{-5\pi u^2}$ is $-12\pi(30\pi^2u^4-19\pi u^2+1)$ times a nonzero factor), and neither does any $L+\kappa X$ with $L$ a local
differential operator and $X=F_Su^2F_S^{-1}$, because $X$ dilates supports outward across the dyadic scales; so the Halmos
angle operator is, for now, the only operator with that commutation, and it is not differential. Fourth, whether the
direction $s_0$ of Observation~\ref{obs:s0}, a purely archimedean object with a frequency drifting like $4/T$, is something you
recognize from the prolate side.

\begin{thebibliography}{99}
\bibitem{Bombieri2000} E.~Bombieri, \emph{Remarks on Weil's quadratic functional in the theory of prime numbers, I}, Rend.\ Lincei Mat.\ Appl.\ \textbf{11} (2000), 183--233.
\bibitem{Burnol1999} J.-F.~Burnol, \emph{The explicit formula in simple terms}, arXiv:math/9810169; \emph{The explicit formula and the conductor operator}, arXiv:math/9902080.
\bibitem{CC2021} A.~Connes and C.~Consani, \emph{Weil positivity and trace formula, the archimedean place}, Selecta Math.\ (N.S.) \textbf{27} (2021), no.~4, Paper No.~77; arXiv:2006.13771.
\bibitem{CCquasi} A.~Connes and C.~Consani, \emph{Quasi-inner functions and local factors}, arXiv:2008.10974.
\bibitem{CCM2023} A.~Connes, C.~Consani and H.~Moscovici, \emph{Zeta zeros and prolate wave operators}, arXiv:2310.18423.
\bibitem{CCM2024} A.~Connes, C.~Consani and H.~Moscovici, arXiv:2403.01247 (2024): the moment problem of the local factors and the semilocal case $S=\{p,\infty\}$.
\bibitem{CM2022} A.~Connes and H.~Moscovici, \emph{The UV prolate spectrum matches the zeros of zeta}, Proc.\ Natl.\ Acad.\ Sci.\ USA (2022).
\bibitem{Connes2026} A.~Connes, arXiv:2602.04022 (2026), survey; \S7.2 (Theorem~7.1) and \S7.4 (semilocal formula).
\bibitem{CvS2025} A.~Connes and W.~D.~van Suijlekom, \emph{Quadratic forms, real zeros and echoes of the spectral action}, arXiv:2511.23257.
\bibitem{Groskin2026} A.~Groskin, \emph{A finite Guinand--Weil dictionary and archimedean tail order for the truncated Weil quadratic form}, arXiv:2607.02828.
\bibitem{Halmos1969} P.~R.~Halmos, \emph{Two subspaces}, Trans.\ Amer.\ Math.\ Soc.\ \textbf{144} (1969), 381--389.
\bibitem{Kim2026} Kim et al., arXiv:2607.24830 (2026): numerical law $\log(1/a)+\log(k-2)+B_0$ for Suzuki's prime-free operator.
\bibitem{Simon} B.~Simon, \emph{Trace Ideals and Their Applications}, 2nd ed., Amer.\ Math.\ Soc., 2005.
\bibitem{Suzuki2023} M.~Suzuki, \emph{Aspects of the screw function corresponding to the Riemann zeta function}, J.~Lond.\ Math.\ Soc.\ (2) \textbf{107} (2023), 1846--1888; arXiv:2206.03682.
\bibitem{Suzuki2026} M.~Suzuki, \emph{Weil's quadratic form via the screw function}, arXiv:2606.09096.
\bibitem{Paper47} D.~A.~Trejo Pizzo, \emph{A dual theory of Weil positivity at finite support: outside-window certificates, non-attainment, and the explicit formula as the dual variable}, preprint, 2026.
\bibitem{Yoshida1992} H.~Yoshida, \emph{On Hermitian forms attached to zeta functions}, in: \emph{Zeta Functions in Geometry}, Adv.\ Stud.\ Pure Math.\ \textbf{21}, Kinokuniya, Tokyo, 1992, 281--325.
\bibitem{Zhu2026} X.~Zhu, \emph{Weil positivity in compact windows: a finite reduction, certified two-sided bounds, and a Landau--Widom decay law}, arXiv:2608.24827.
\end{thebibliography}

\end{document}
